% Requires:
% \usepackage{amsmath}
% \usetikzlibrary{calc}
% \usetikzlibrary{intersections}
% \usetikzlibrary{patterns.meta}
% \usetikzlibrary{shapes.misc}
% \providecommand{\glsxtrshort}[1]{#1}
% \providecommand{\Glsxtrlong}[1]{#1}
\begin{tikzpicture}[
    scale=0.75,
    ray/.style={->,thick,draw=green!80!black,opacity=0.75},
    invalid ray/.style={ray,dashed},
    wall/.style={draw=red, thick},
    point/.style={draw,cross out,thick,inner sep=2pt},
    ground/.style={fill,draw=none,minimum width=0.3,minimum height=0.6},
    intersection-point/.style={draw=blue,fill=white,fill opacity=0.5,circle,thick,inner sep=2pt},
    previous-intersection-point/.style={intersection-point,draw=blue!50!white,opacity=#1}
  ]
  % Macro to draw the static background elements (walls, ground, tx, rx)
  \def\drawbackground{
    % Wall 1 (Angled)
    \draw[ground,pattern={Lines[angle=0]},rotate around={45:(0,0)}] (0,0) rectangle ({3.5/cos(45)}, .1);
    \draw[wall] (0,0) -- (3.5, 3.5);

    % Wall 2 (Vertical)
    \draw[ground,pattern={Lines[angle=45]}] (5,.5) rectangle (5.1, 4);
    \draw[wall] (5,.5) -- (5, 4);

    % Tx and Rx points
    \node[point,label=below:{Transmitter}] (BS) at (2,-1) {};
    \node[point,label=above:{Receiver}] (UE) at (2,4) {};
  }

  % ---------------------------------------------------------
  % Iteration 1: C = 2.09
  % ---------------------------------------------------------
  \drawbackground

  % Approximate Specular Points
  \coordinate (S1) at (1.5, 1.5);
  \coordinate (S2) at (5, 2.0);

  % Ray
  \draw[invalid ray] (BS) -- (S1) -- (S2) -- (UE);

  % Intersection Points
  \node[intersection-point] at (S1) {};
  \node[intersection-point] at (S2) {};

  % Cost Label
  \node[anchor=west] at (5.5, 2.5) {$C = 2.09$};

\end{tikzpicture}
